% --- Archon physics formalization source begin ---
% archon:physics
% archon:covers PhyXMiniProblems/problem_phyx_mini_0299.lean
% archon:source-report reports/phyx_mini/problem_phyx_mini_0299.source.json

\chapter{Physics problem phyx\_mini\_0299}
\label{ch:PhyXMiniProblems_problem_phyx_mini_0299}

\paragraph{Problem source.}
\#\# Physical scenario

In figure, an aluminum wire, of length \$L\_1 = 60.0\$ cm, cross-sectional area \$1.00 \textbackslash{}times 10\textasciicircum{}\{-2\}\$ cm\$\textasciicircum{}2\$, and density \$2.60\$ g/cm\$\textasciicircum{}3\$, is joined to a steel wire, of density \$7.80\$ g/cm\$\textasciicircum{}3\$ and the same cross-sectional area. The compound wire, loaded with a block of mass \$m = 10.0\$ kg, is arranged so that the distance \$L\_2\$ from the joint to the supporting pulley is \$86.6\$ cm. Transverse waves are set up on the wire by an external source of variable frequency; a node is located at the pulley.

\#\# Question

Find the lowest frequency that generates a standing wave having the joint as one of the nodes.

\#\# Answer choices

- A: A: 305 Hz

- B: B: 312 Hz

- C: C: 318 Hz

- D: D: 324 Hz

\#\# Figure caption (auxiliary; use the image as primary evidence)

The image depicts a mechanical setup involving a pulley and a block system. Here's a detailed description:

- **Pulley System**: A pulley is attached to the edge of a horizontal, rectangular beam.
- **Rope/Cable**: Two sections of a rope or cable are shown. The first section, labeled "Aluminum," covers a distance labeled \textbackslash{}( L\_1 \textbackslash{}) and rests atop the beam. The second section, labeled "Steel," covers a distance labeled \textbackslash{}( L\_2 \textbackslash{}) and also rests on the beam before passing over the pulley.
- **Block**: A block labeled with mass \textbackslash{}( m \textbackslash{}) is suspended from the end of the steel cable, hanging vertically.
- **Materials \& Labels**: The beam and the inner wall of the system where the pulley is attached is colored brown, while the distinct sections of materials (aluminum and steel) indicate a practical or theoretical examination of differing properties (possibly tension or weight effects), with lengths \textbackslash{}( L\_1 \textbackslash{}) and \textbackslash{}( L\_2 \textbackslash{}).
  
The setup likely illustrates concepts related to mechanics or physics, such as tension, material properties, or mechanical advantage in a pulley system.

\#\# Dataset metadata

Resonance and Harmonics; Multi-Formula Reasoning, Physical Model Grounding Reasoning

\paragraph{Recorded answer/context.}
D: D: 324 Hz

\paragraph{Figure/image path.}
/root/proposal\_for\_physic/science-mango/phyx\_mini\_run/phyx\_data/test\_image/299.png

\paragraph{Formalization target.}
create a compiling Lean file with sorry bodies at `PhyXMiniProblems/problem\_phyx\_mini\_0299.lean`. The Lean declarations must preserve the physical quantities, dimensions or dimensional roles, figure labels, governing-law hypotheses, and final relation expressed by this problem.
Use Mathlib/Physlib names found through LeanExplore where available. If a physics API is missing, introduce faithful local abstractions rather than scalar placeholder aliases.

\begin{theorem}[Physics formalization target]
\label{thm:physics:phyx_mini_0299:target}
\lean{PhyXMiniProblems.ProblemPhyXMini0299.problem_phyx_mini_0299}
\uses{def:physics:phyx-mini-0299:phyxminiproblems-problemphyxmini0299-frequencyinhertz, def:physics:phyx-mini-0299:phyxminiproblems-problemphyxmini0299-compoundwiresetup, def:physics:phyx-mini-0299:phyxminiproblems-problemphyxmini0299-matchesproblemdata, def:physics:phyx-mini-0299:phyxminiproblems-problemphyxmini0299-usesstandardterrestrialgravity, def:physics:phyx-mini-0299:phyxminiproblems-problemphyxmini0299-matchesprimaryfigure, def:physics:phyx-mini-0299:phyxminiproblems-problemphyxmini0299-haspositivephysicalparameters, def:physics:phyx-mini-0299:phyxminiproblems-problemphyxmini0299-satisfiescompoundwirelaws, def:physics:phyx-mini-0299:phyxminiproblems-problemphyxmini0299-modepairmatcheswholehertz, def:physics:phyx-mini-0299:phyxminiproblems-problemphyxmini0299-islowestjointnodefrequencyatwholehertz, def:physics:phyx-mini-0299:phyxminiproblems-problemphyxmini0299-recordedanswerchoice}
The assigned autoformalize agent should translate this physics problem into Lean declarations in the covered file, with theorem and lemma proof bodies written as `by sorry`.
\end{theorem}
\begin{proof}
This is an autoformalization task, not a proof task. Produce statements that can later be proved by the physics prover without weakening the physical model.
\end{proof}

% --- Archon declaration topology begin ---
\section{Declaration topology}
\begin{definition}[Length Quantity]
  \label{def:physics:phyx-mini-0299:phyxminiproblems-problemphyxmini0299-lengthquantity}
  \lean{PhyXMiniProblems.ProblemPhyXMini0299.LengthQuantity}
  A nonnegative physical length
\end{definition}
\begin{proof}
  This is the declared model definition of Length Quantity.
\end{proof}
\begin{definition}[Area Quantity]
  \label{def:physics:phyx-mini-0299:phyxminiproblems-problemphyxmini0299-areaquantity}
  \lean{PhyXMiniProblems.ProblemPhyXMini0299.AreaQuantity}
  A nonnegative cross-sectional area
\end{definition}
\begin{proof}
  This is the declared model definition of Area Quantity.
\end{proof}
\begin{definition}[Volume Mass Density Quantity]
  \label{def:physics:phyx-mini-0299:phyxminiproblems-problemphyxmini0299-volumemassdensityquantity}
  \lean{PhyXMiniProblems.ProblemPhyXMini0299.VolumeMassDensityQuantity}
  A nonnegative mass per unit volume
\end{definition}
\begin{proof}
  This is the declared model definition of Volume Mass Density Quantity.
\end{proof}
\begin{definition}[Linear Mass Density Quantity]
  \label{def:physics:phyx-mini-0299:phyxminiproblems-problemphyxmini0299-linearmassdensityquantity}
  \lean{PhyXMiniProblems.ProblemPhyXMini0299.LinearMassDensityQuantity}
  A nonnegative mass per unit length
\end{definition}
\begin{proof}
  This is the declared model definition of Linear Mass Density Quantity.
\end{proof}
\begin{definition}[Mass Quantity]
  \label{def:physics:phyx-mini-0299:phyxminiproblems-problemphyxmini0299-massquantity}
  \lean{PhyXMiniProblems.ProblemPhyXMini0299.MassQuantity}
  A nonnegative physical mass
\end{definition}
\begin{proof}
  This is the declared model definition of Mass Quantity.
\end{proof}
\begin{definition}[Acceleration Quantity]
  \label{def:physics:phyx-mini-0299:phyxminiproblems-problemphyxmini0299-accelerationquantity}
  \lean{PhyXMiniProblems.ProblemPhyXMini0299.AccelerationQuantity}
  A nonnegative acceleration magnitude
\end{definition}
\begin{proof}
  This is the declared model definition of Acceleration Quantity.
\end{proof}
\begin{definition}[Tension Quantity]
  \label{def:physics:phyx-mini-0299:phyxminiproblems-problemphyxmini0299-tensionquantity}
  \lean{PhyXMiniProblems.ProblemPhyXMini0299.TensionQuantity}
  A nonnegative tensile force
\end{definition}
\begin{proof}
  This is the declared model definition of Tension Quantity.
\end{proof}
\begin{definition}[Frequency Quantity]
  \label{def:physics:phyx-mini-0299:phyxminiproblems-problemphyxmini0299-frequencyquantity}
  \lean{PhyXMiniProblems.ProblemPhyXMini0299.FrequencyQuantity}
  An ordinary cyclic frequency, with inverse-time dimension
\end{definition}
\begin{proof}
  This is the declared model definition of Frequency Quantity.
\end{proof}
\begin{definition}[Speed Quantity]
  \label{def:physics:phyx-mini-0299:phyxminiproblems-problemphyxmini0299-speedquantity}
  \lean{PhyXMiniProblems.ProblemPhyXMini0299.SpeedQuantity}
  A nonnegative transverse-wave propagation speed
\end{definition}
\begin{proof}
  This is the declared model definition of Speed Quantity.
\end{proof}
\begin{definition}[quantity Readout]
  \label{def:physics:phyx-mini-0299:phyxminiproblems-problemphyxmini0299-quantityreadout}
  \lean{PhyXMiniProblems.ProblemPhyXMini0299.quantityReadout}
  Read a nonnegative dimensionful quantity in a coherent unit system
\end{definition}
\begin{proof}
  This is the declared model definition of quantity Readout.
\end{proof}
\begin{definition}[length In Meters]
  \label{def:physics:phyx-mini-0299:phyxminiproblems-problemphyxmini0299-lengthinmeters}
  \lean{PhyXMiniProblems.ProblemPhyXMini0299.lengthInMeters}
  \uses{def:physics:phyx-mini-0299:phyxminiproblems-problemphyxmini0299-lengthquantity, def:physics:phyx-mini-0299:phyxminiproblems-problemphyxmini0299-quantityreadout}
  SI metre readout of a length
\end{definition}
\begin{proof}
  This is the declared model definition of length In Meters.
\end{proof}
\begin{definition}[area In Square Meters]
  \label{def:physics:phyx-mini-0299:phyxminiproblems-problemphyxmini0299-areainsquaremeters}
  \lean{PhyXMiniProblems.ProblemPhyXMini0299.areaInSquareMeters}
  \uses{def:physics:phyx-mini-0299:phyxminiproblems-problemphyxmini0299-areaquantity, def:physics:phyx-mini-0299:phyxminiproblems-problemphyxmini0299-quantityreadout}
  SI square-metre readout of an area
\end{definition}
\begin{proof}
  This is the declared model definition of area In Square Meters.
\end{proof}
\begin{definition}[volume Density In Kilograms Per Cubic Meter]
  \label{def:physics:phyx-mini-0299:phyxminiproblems-problemphyxmini0299-volumedensityinkilogramspercubicmeter}
  \lean{PhyXMiniProblems.ProblemPhyXMini0299.volumeDensityInKilogramsPerCubicMeter}
  \uses{def:physics:phyx-mini-0299:phyxminiproblems-problemphyxmini0299-volumemassdensityquantity, def:physics:phyx-mini-0299:phyxminiproblems-problemphyxmini0299-quantityreadout}
  SI kilogram-per-cubic-metre readout of a volume mass density
\end{definition}
\begin{proof}
  This is the declared model definition of volume Density In Kilograms Per Cubic Meter.
\end{proof}
\begin{definition}[linear Density In Kilograms Per Meter]
  \label{def:physics:phyx-mini-0299:phyxminiproblems-problemphyxmini0299-lineardensityinkilogramspermeter}
  \lean{PhyXMiniProblems.ProblemPhyXMini0299.linearDensityInKilogramsPerMeter}
  \uses{def:physics:phyx-mini-0299:phyxminiproblems-problemphyxmini0299-linearmassdensityquantity, def:physics:phyx-mini-0299:phyxminiproblems-problemphyxmini0299-quantityreadout}
  SI kilogram-per-metre readout of a linear mass density
\end{definition}
\begin{proof}
  This is the declared model definition of linear Density In Kilograms Per Meter.
\end{proof}
\begin{definition}[mass In Kilograms]
  \label{def:physics:phyx-mini-0299:phyxminiproblems-problemphyxmini0299-massinkilograms}
  \lean{PhyXMiniProblems.ProblemPhyXMini0299.massInKilograms}
  \uses{def:physics:phyx-mini-0299:phyxminiproblems-problemphyxmini0299-massquantity, def:physics:phyx-mini-0299:phyxminiproblems-problemphyxmini0299-quantityreadout}
  Kilogram readout of a mass
\end{definition}
\begin{proof}
  This is the declared model definition of mass In Kilograms.
\end{proof}
\begin{definition}[acceleration In Meters Per Second Squared]
  \label{def:physics:phyx-mini-0299:phyxminiproblems-problemphyxmini0299-accelerationinmeterspersecondsquared}
  \lean{PhyXMiniProblems.ProblemPhyXMini0299.accelerationInMetersPerSecondSquared}
  \uses{def:physics:phyx-mini-0299:phyxminiproblems-problemphyxmini0299-accelerationquantity, def:physics:phyx-mini-0299:phyxminiproblems-problemphyxmini0299-quantityreadout}
  SI metre-per-second-squared readout of an acceleration
\end{definition}
\begin{proof}
  This is the declared model definition of acceleration In Meters Per Second Squared.
\end{proof}
\begin{definition}[tension In Newtons]
  \label{def:physics:phyx-mini-0299:phyxminiproblems-problemphyxmini0299-tensioninnewtons}
  \lean{PhyXMiniProblems.ProblemPhyXMini0299.tensionInNewtons}
  \uses{def:physics:phyx-mini-0299:phyxminiproblems-problemphyxmini0299-tensionquantity, def:physics:phyx-mini-0299:phyxminiproblems-problemphyxmini0299-quantityreadout}
  SI newton readout of a tension
\end{definition}
\begin{proof}
  This is the declared model definition of tension In Newtons.
\end{proof}
\begin{definition}[frequency In Hertz]
  \label{def:physics:phyx-mini-0299:phyxminiproblems-problemphyxmini0299-frequencyinhertz}
  \lean{PhyXMiniProblems.ProblemPhyXMini0299.frequencyInHertz}
  \uses{def:physics:phyx-mini-0299:phyxminiproblems-problemphyxmini0299-frequencyquantity, def:physics:phyx-mini-0299:phyxminiproblems-problemphyxmini0299-quantityreadout}
  SI hertz readout of an ordinary frequency
\end{definition}
\begin{proof}
  This is the declared model definition of frequency In Hertz.
\end{proof}
\begin{definition}[speed In Meters Per Second]
  \label{def:physics:phyx-mini-0299:phyxminiproblems-problemphyxmini0299-speedinmeterspersecond}
  \lean{PhyXMiniProblems.ProblemPhyXMini0299.speedInMetersPerSecond}
  \uses{def:physics:phyx-mini-0299:phyxminiproblems-problemphyxmini0299-speedquantity, def:physics:phyx-mini-0299:phyxminiproblems-problemphyxmini0299-quantityreadout}
  SI metre-per-second readout of a speed
\end{definition}
\begin{proof}
  This is the declared model definition of speed In Meters Per Second.
\end{proof}
\begin{definition}[Wire Segment]
  \label{def:physics:phyx-mini-0299:phyxminiproblems-problemphyxmini0299-wiresegment}
  \lean{PhyXMiniProblems.ProblemPhyXMini0299.WireSegment}
  The two horizontal portions of the compound wire
\end{definition}
\begin{proof}
  This is the declared model definition of Wire Segment.
\end{proof}
\begin{definition}[Wire Material]
  \label{def:physics:phyx-mini-0299:phyxminiproblems-problemphyxmini0299-wirematerial}
  \lean{PhyXMiniProblems.ProblemPhyXMini0299.WireMaterial}
  Material names printed below the two horizontal wire segments
\end{definition}
\begin{proof}
  This is the declared model definition of Wire Material.
\end{proof}
\begin{definition}[Figure Length Label]
  \label{def:physics:phyx-mini-0299:phyxminiproblems-problemphyxmini0299-figurelengthlabel}
  \lean{PhyXMiniProblems.ProblemPhyXMini0299.FigureLengthLabel}
  Length symbols printed above the two horizontal segments
\end{definition}
\begin{proof}
  This is the declared model definition of Figure Length Label.
\end{proof}
\begin{definition}[Figure Point]
  \label{def:physics:phyx-mini-0299:phyxminiproblems-problemphyxmini0299-figurepoint}
  \lean{PhyXMiniProblems.ProblemPhyXMini0299.FigurePoint}
  Distinguished locations or objects visible in the primary image
\end{definition}
\begin{proof}
  This is the declared model definition of Figure Point.
\end{proof}
\begin{definition}[Figure Connection]
  \label{def:physics:phyx-mini-0299:phyxminiproblems-problemphyxmini0299-figureconnection}
  \lean{PhyXMiniProblems.ProblemPhyXMini0299.FigureConnection}
  Mechanical role of each distinguished image location
\end{definition}
\begin{proof}
  This is the declared model definition of Figure Connection.
\end{proof}
\begin{definition}[Wire Route]
  \label{def:physics:phyx-mini-0299:phyxminiproblems-problemphyxmini0299-wireroute}
  \lean{PhyXMiniProblems.ProblemPhyXMini0299.WireRoute}
  The route taken by the compound wire in the primary image
\end{definition}
\begin{proof}
  This is the declared model definition of Wire Route.
\end{proof}
\begin{definition}[Drive Kind]
  \label{def:physics:phyx-mini-0299:phyxminiproblems-problemphyxmini0299-drivekind}
  \lean{PhyXMiniProblems.ProblemPhyXMini0299.DriveKind}
  The drive described in the problem text
\end{definition}
\begin{proof}
  This is the declared model definition of Drive Kind.
\end{proof}
\begin{definition}[Transverse Boundary Condition]
  \label{def:physics:phyx-mini-0299:phyxminiproblems-problemphyxmini0299-transverseboundarycondition}
  \lean{PhyXMiniProblems.ProblemPhyXMini0299.TransverseBoundaryCondition}
  Transverse-displacement behavior at a supported point
\end{definition}
\begin{proof}
  This is the declared model definition of Transverse Boundary Condition.
\end{proof}
\begin{definition}[Source Figure]
  \label{def:physics:phyx-mini-0299:phyxminiproblems-problemphyxmini0299-sourcefigure}
  \lean{PhyXMiniProblems.ProblemPhyXMini0299.SourceFigure}
  \uses{def:physics:phyx-mini-0299:phyxminiproblems-problemphyxmini0299-wiresegment, def:physics:phyx-mini-0299:phyxminiproblems-problemphyxmini0299-wirematerial, def:physics:phyx-mini-0299:phyxminiproblems-problemphyxmini0299-figurelengthlabel, def:physics:phyx-mini-0299:phyxminiproblems-problemphyxmini0299-figurepoint, def:physics:phyx-mini-0299:phyxminiproblems-problemphyxmini0299-figureconnection, def:physics:phyx-mini-0299:phyxminiproblems-problemphyxmini0299-wireroute}
  Structured transcription of labels and connections in the primary image
\end{definition}
\begin{proof}
  This is the declared model definition of Source Figure.
\end{proof}
\begin{definition}[Compound Wire Setup]
  \label{def:physics:phyx-mini-0299:phyxminiproblems-problemphyxmini0299-compoundwiresetup}
  \lean{PhyXMiniProblems.ProblemPhyXMini0299.CompoundWireSetup}
  \uses{def:physics:phyx-mini-0299:phyxminiproblems-problemphyxmini0299-lengthquantity, def:physics:phyx-mini-0299:phyxminiproblems-problemphyxmini0299-areaquantity, def:physics:phyx-mini-0299:phyxminiproblems-problemphyxmini0299-volumemassdensityquantity, def:physics:phyx-mini-0299:phyxminiproblems-problemphyxmini0299-linearmassdensityquantity, def:physics:phyx-mini-0299:phyxminiproblems-problemphyxmini0299-massquantity, def:physics:phyx-mini-0299:phyxminiproblems-problemphyxmini0299-accelerationquantity, def:physics:phyx-mini-0299:phyxminiproblems-problemphyxmini0299-tensionquantity, def:physics:phyx-mini-0299:phyxminiproblems-problemphyxmini0299-frequencyquantity, def:physics:phyx-mini-0299:phyxminiproblems-problemphyxmini0299-speedquantity, def:physics:phyx-mini-0299:phyxminiproblems-problemphyxmini0299-wiresegment, def:physics:phyx-mini-0299:phyxminiproblems-problemphyxmini0299-wirematerial, def:physics:phyx-mini-0299:phyxminiproblems-problemphyxmini0299-drivekind, def:physics:phyx-mini-0299:phyxminiproblems-problemphyxmini0299-transverseboundarycondition, def:physics:phyx-mini-0299:phyxminiproblems-problemphyxmini0299-sourcefigure}
  Independent physical quantities for the compound-wire apparatus. nodeToNodeModeFrequency segment n is the physical frequency of the local nth node-to-node standing-wave mode of that segment. Its generic relation to length and wave speed is imposed below. No common mode, lowest mode, or answer value is selected in this structure
\end{definition}
\begin{proof}
  This is the declared model definition of Compound Wire Setup.
\end{proof}
\begin{definition}[Matches Problem Data]
  \label{def:physics:phyx-mini-0299:phyxminiproblems-problemphyxmini0299-matchesproblemdata}
  \lean{PhyXMiniProblems.ProblemPhyXMini0299.MatchesProblemData}
  \uses{def:physics:phyx-mini-0299:phyxminiproblems-problemphyxmini0299-lengthinmeters, def:physics:phyx-mini-0299:phyxminiproblems-problemphyxmini0299-areainsquaremeters, def:physics:phyx-mini-0299:phyxminiproblems-problemphyxmini0299-volumedensityinkilogramspercubicmeter, def:physics:phyx-mini-0299:phyxminiproblems-problemphyxmini0299-massinkilograms, def:physics:phyx-mini-0299:phyxminiproblems-problemphyxmini0299-compoundwiresetup}
  The numerical data printed in the problem, converted to exact SI readouts of the printed decimals: 60.0 cm = 3/5 m and 86.6 cm = 433/500 m; 1.00 × 10⁻² cm² = 10⁻⁶ m²; 2.60 g/cm³ = 2600 kg/m³ and 7.80 g/cm³ = 7800 kg/m³; m = 10.0 kg. The source is variable-frequency and transverse, and the problem explicitly states that the pulley is a node. No condition at the material joint occurs in this data structure
\end{definition}
\begin{proof}
  This is the declared model definition of Matches Problem Data.
\end{proof}
\begin{definition}[Uses Standard Terrestrial Gravity]
  \label{def:physics:phyx-mini-0299:phyxminiproblems-problemphyxmini0299-usesstandardterrestrialgravity}
  \lean{PhyXMiniProblems.ProblemPhyXMini0299.UsesStandardTerrestrialGravity}
  \uses{def:physics:phyx-mini-0299:phyxminiproblems-problemphyxmini0299-accelerationinmeterspersecondsquared, def:physics:phyx-mini-0299:phyxminiproblems-problemphyxmini0299-compoundwiresetup}
  The usual textbook value g = 9.80 m/s² used with the hanging 10.0 kg block. This is an environmental calibration, not a frequency conclusion
\end{definition}
\begin{proof}
  This is the declared model definition of Uses Standard Terrestrial Gravity.
\end{proof}
\begin{definition}[Matches Primary Figure]
  \label{def:physics:phyx-mini-0299:phyxminiproblems-problemphyxmini0299-matchesprimaryfigure}
  \lean{PhyXMiniProblems.ProblemPhyXMini0299.MatchesPrimaryFigure}
  \uses{def:physics:phyx-mini-0299:phyxminiproblems-problemphyxmini0299-compoundwiresetup}
  Primary-image readout: a left fixed support, an aluminum L₁ segment, the material joint, a steel L₂ segment, a pulley, and a hanging block labelled m, in that route order
\end{definition}
\begin{proof}
  This is the declared model definition of Matches Primary Figure.
\end{proof}
\begin{definition}[Has Positive Physical Parameters]
  \label{def:physics:phyx-mini-0299:phyxminiproblems-problemphyxmini0299-haspositivephysicalparameters}
  \lean{PhyXMiniProblems.ProblemPhyXMini0299.HasPositivePhysicalParameters}
  \uses{def:physics:phyx-mini-0299:phyxminiproblems-problemphyxmini0299-lengthinmeters, def:physics:phyx-mini-0299:phyxminiproblems-problemphyxmini0299-areainsquaremeters, def:physics:phyx-mini-0299:phyxminiproblems-problemphyxmini0299-volumedensityinkilogramspercubicmeter, def:physics:phyx-mini-0299:phyxminiproblems-problemphyxmini0299-lineardensityinkilogramspermeter, def:physics:phyx-mini-0299:phyxminiproblems-problemphyxmini0299-massinkilograms, def:physics:phyx-mini-0299:phyxminiproblems-problemphyxmini0299-accelerationinmeterspersecondsquared, def:physics:phyx-mini-0299:phyxminiproblems-problemphyxmini0299-tensioninnewtons, def:physics:phyx-mini-0299:phyxminiproblems-problemphyxmini0299-frequencyinhertz, def:physics:phyx-mini-0299:phyxminiproblems-problemphyxmini0299-speedinmeterspersecond, def:physics:phyx-mini-0299:phyxminiproblems-problemphyxmini0299-compoundwiresetup}
  Positivity and nondegeneracy conditions for the taut-wire model
\end{definition}
\begin{proof}
  This is the declared model definition of Has Positive Physical Parameters.
\end{proof}
\begin{definition}[segment Material]
  \label{def:physics:phyx-mini-0299:phyxminiproblems-problemphyxmini0299-segmentmaterial}
  \lean{PhyXMiniProblems.ProblemPhyXMini0299.segmentMaterial}
  \uses{def:physics:phyx-mini-0299:phyxminiproblems-problemphyxmini0299-wiresegment, def:physics:phyx-mini-0299:phyxminiproblems-problemphyxmini0299-wirematerial}
  Material associated to each segment independently of the source figure
\end{definition}
\begin{proof}
  This is the declared model definition of segment Material.
\end{proof}
\begin{definition}[Satisfies Compound Wire Laws]
  \label{def:physics:phyx-mini-0299:phyxminiproblems-problemphyxmini0299-satisfiescompoundwirelaws}
  \lean{PhyXMiniProblems.ProblemPhyXMini0299.SatisfiesCompoundWireLaws}
  \uses{def:physics:phyx-mini-0299:phyxminiproblems-problemphyxmini0299-quantityreadout, def:physics:phyx-mini-0299:phyxminiproblems-problemphyxmini0299-compoundwiresetup, def:physics:phyx-mini-0299:phyxminiproblems-problemphyxmini0299-segmentmaterial}
  The standard physical laws used in the calculation, expressed in every coherent unit system: linear density is volume density times common cross-sectional area; an ideal pulley and a static hanging block give each segment tension mg; a taut string obeys μ v² = T; the nth node-to-node mode obeys 2 L fₙ = n v. These laws describe the full modal spectra. They do not assert that a mode of one material coincides with a mode of the other, do not select a lowest common mode, and contain neither 324 nor an answer label
\end{definition}
\begin{proof}
  This is the declared model definition of Satisfies Compound Wire Laws.
\end{proof}
\begin{definition}[Mode Pair Matches Whole Hertz]
  \label{def:physics:phyx-mini-0299:phyxminiproblems-problemphyxmini0299-modepairmatcheswholehertz}
  \lean{PhyXMiniProblems.ProblemPhyXMini0299.ModePairMatchesWholeHertz}
  \uses{def:physics:phyx-mini-0299:phyxminiproblems-problemphyxmini0299-frequencyinhertz, def:physics:phyx-mini-0299:phyxminiproblems-problemphyxmini0299-compoundwiresetup}
  A positive aluminum/steel harmonic pair agrees with a whole-hertz drive readout when each local node-to-node mode lies within half a hertz of it
\end{definition}
\begin{proof}
  This is the declared model definition of Mode Pair Matches Whole Hertz.
\end{proof}
\begin{definition}[Is Joint Node Standing Wave At Whole Hertz]
  \label{def:physics:phyx-mini-0299:phyxminiproblems-problemphyxmini0299-isjointnodestandingwaveatwholehertz}
  \lean{PhyXMiniProblems.ProblemPhyXMini0299.IsJointNodeStandingWaveAtWholeHertz}
  \uses{def:physics:phyx-mini-0299:phyxminiproblems-problemphyxmini0299-compoundwiresetup, def:physics:phyx-mini-0299:phyxminiproblems-problemphyxmini0299-modepairmatcheswholehertz}
  At the precision of the integer-hertz answers, a common joint-node standing wave is a shared drive readout matching a node-to-node mode of each segment. The material joint is therefore an endpoint node of both local modes
\end{definition}
\begin{proof}
  This is the declared model definition of Is Joint Node Standing Wave At Whole Hertz.
\end{proof}
\begin{definition}[Is Lowest Joint Node Frequency At Whole Hertz]
  \label{def:physics:phyx-mini-0299:phyxminiproblems-problemphyxmini0299-islowestjointnodefrequencyatwholehertz}
  \lean{PhyXMiniProblems.ProblemPhyXMini0299.IsLowestJointNodeFrequencyAtWholeHertz}
  \uses{def:physics:phyx-mini-0299:phyxminiproblems-problemphyxmini0299-compoundwiresetup, def:physics:phyx-mini-0299:phyxminiproblems-problemphyxmini0299-isjointnodestandingwaveatwholehertz}
  A common joint-node displayed frequency with no lower positive one
\end{definition}
\begin{proof}
  This is the declared model definition of Is Lowest Joint Node Frequency At Whole Hertz.
\end{proof}
\begin{definition}[Answer Choice]
  \label{def:physics:phyx-mini-0299:phyxminiproblems-problemphyxmini0299-answerchoice}
  \lean{PhyXMiniProblems.ProblemPhyXMini0299.AnswerChoice}
  Labels of the four displayed answers
\end{definition}
\begin{proof}
  This is the declared model definition of Answer Choice.
\end{proof}
\begin{definition}[frequency In Hertz]
  \label{def:physics:phyx-mini-0299:phyxminiproblems-problemphyxmini0299-answerchoice-frequencyinhertz}
  \lean{PhyXMiniProblems.ProblemPhyXMini0299.AnswerChoice.frequencyInHertz}
  \uses{def:physics:phyx-mini-0299:phyxminiproblems-problemphyxmini0299-frequencyinhertz, def:physics:phyx-mini-0299:phyxminiproblems-problemphyxmini0299-answerchoice}
  Integer-hertz frequency printed beside an answer label
\end{definition}
\begin{proof}
  This is the declared model definition of frequency In Hertz.
\end{proof}
\begin{definition}[recorded Answer Choice]
  \label{def:physics:phyx-mini-0299:phyxminiproblems-problemphyxmini0299-recordedanswerchoice}
  \lean{PhyXMiniProblems.ProblemPhyXMini0299.recordedAnswerChoice}
  \uses{def:physics:phyx-mini-0299:phyxminiproblems-problemphyxmini0299-answerchoice}
  The answer label recorded in the supplied dataset metadata
\end{definition}
\begin{proof}
  This is the declared model definition of recorded Answer Choice.
\end{proof}
\begin{lemma}[second aluminum and fifth steel modes match choice D]
  \label{lem:physics:phyx-mini-0299:phyxminiproblems-problemphyxmini0299-second-aluminum-and-fifth-steel-modes-match-choice-d}
  \lean{PhyXMiniProblems.ProblemPhyXMini0299.second_aluminum_and_fifth_steel_modes_match_choice_D}
  \uses{def:physics:phyx-mini-0299:phyxminiproblems-problemphyxmini0299-frequencyinhertz, def:physics:phyx-mini-0299:phyxminiproblems-problemphyxmini0299-compoundwiresetup, def:physics:phyx-mini-0299:phyxminiproblems-problemphyxmini0299-matchesproblemdata, def:physics:phyx-mini-0299:phyxminiproblems-problemphyxmini0299-usesstandardterrestrialgravity, def:physics:phyx-mini-0299:phyxminiproblems-problemphyxmini0299-haspositivephysicalparameters, def:physics:phyx-mini-0299:phyxminiproblems-problemphyxmini0299-satisfiescompoundwirelaws, def:physics:phyx-mini-0299:phyxminiproblems-problemphyxmini0299-modepairmatcheswholehertz, def:physics:phyx-mini-0299:phyxminiproblems-problemphyxmini0299-answerchoice}
  The second aluminum mode and fifth steel mode both agree with the displayed frequency 324 Hz at whole-hertz precision
\end{lemma}
\begin{proof}
  The conclusion follows from the governing relations and figure readouts named in the statement of second aluminum and fifth steel modes match choice D.
\end{proof}
\begin{lemma}[every joint node frequency is at least choice D]
  \label{lem:physics:phyx-mini-0299:phyxminiproblems-problemphyxmini0299-every-joint-node-frequency-is-at-least-choice-d}
  \lean{PhyXMiniProblems.ProblemPhyXMini0299.every_joint_node_frequency_is_at_least_choice_D}
  \uses{def:physics:phyx-mini-0299:phyxminiproblems-problemphyxmini0299-frequencyinhertz, def:physics:phyx-mini-0299:phyxminiproblems-problemphyxmini0299-compoundwiresetup, def:physics:phyx-mini-0299:phyxminiproblems-problemphyxmini0299-matchesproblemdata, def:physics:phyx-mini-0299:phyxminiproblems-problemphyxmini0299-usesstandardterrestrialgravity, def:physics:phyx-mini-0299:phyxminiproblems-problemphyxmini0299-haspositivephysicalparameters, def:physics:phyx-mini-0299:phyxminiproblems-problemphyxmini0299-satisfiescompoundwirelaws, def:physics:phyx-mini-0299:phyxminiproblems-problemphyxmini0299-isjointnodestandingwaveatwholehertz, def:physics:phyx-mini-0299:phyxminiproblems-problemphyxmini0299-answerchoice}
  No lower positive whole-hertz readout simultaneously matches node-to-node modes of both materials
\end{lemma}
\begin{proof}
  The conclusion follows from the governing relations and figure readouts named in the statement of every joint node frequency is at least choice D.
\end{proof}
% --- Archon declaration topology end ---
% --- Archon physics formalization source end ---
